\chapter{System Design and Implementation}
\label{ch:system_design}
\label{ch:reliability}

\section{End-to-End Architecture}
\label{sec:design_architecture}

The framework follows a hub-and-spoke architecture. Edge agents run beside
monitored services and perform metric collection, preprocessing, LSTM-AE
inference, threshold decisions, and local training. The central coordinator
manages federated rounds, persists model versions, maintains the causal graph,
and enriches anomaly alerts with RCA results.

\begin{figure}[H]
\centering
\begin{tikzpicture}[
    node distance=1.2cm and 2.0cm,
    box/.style={draw, rounded corners, minimum width=2.8cm, minimum height=0.9cm,
                align=center, font=\small},
    edge_box/.style={box, fill=blue!8},
    coord_box/.style={box, fill=orange!10},
    data_box/.style={box, fill=green!8},
    arrow/.style={-{Stealth[length=2.5mm]}, thick},
]
\node[edge_box] (e1) {Edge Agent 1\\[-2pt]\scriptsize LSTM-AE + Q-Thresh};
\node[edge_box, below=0.6cm of e1] (e2) {Edge Agent 2\\[-2pt]\scriptsize LSTM-AE + Q-Thresh};
\node[edge_box, below=0.6cm of e2] (e3) {Edge Agent $K$\\[-2pt]\scriptsize LSTM-AE + Q-Thresh};

\node[coord_box, right=3.5cm of e2] (coord) {Central\\Coordinator};
\node[coord_box, above=0.7cm of coord] (fl) {FL Aggregator\\[-2pt]\scriptsize FedAvg + DP};
\node[coord_box, below=0.7cm of coord] (rca) {RCA Engine\\[-2pt]\scriptsize Causal Graph};

\node[data_box, right=2.5cm of coord] (db) {Model Registry\\[-2pt]\scriptsize SQLite / WAL};
\node[data_box, right=2.5cm of fl] (graph) {Causal Graph\\[-2pt]\scriptsize NetworkX};
\node[data_box, below=0.7cm of rca] (traces) {Trace Collector\\[-2pt]\scriptsize OpenTelemetry};

\draw[arrow] (e1.east) -- node[above, font=\scriptsize, sloped] {model updates} (fl.west);
\draw[arrow] (e2.east) -- (coord.west);
\draw[arrow] (e3.east) -- node[below, font=\scriptsize, sloped] {anomaly events} (rca.west);
\draw[arrow] (fl) -- (coord);
\draw[arrow] (rca) -- (coord);
\draw[arrow] (coord) -- (db);
\draw[arrow] (fl) -- (graph);
\draw[arrow] (traces) -- (rca);
\draw[arrow, dashed] (coord.west) -- node[above, font=\scriptsize, sloped] {global model} (e2.east);
\end{tikzpicture}
\caption{High-level architecture of the distributed anomaly detection framework}
\label{fig:architecture_overview}
\end{figure}

A complete operational cycle proceeds as follows:

\begin{enumerate}
    \item the edge agent collects and normalizes telemetry into a sliding window;
    \item the local LSTM-AE produces a reconstruction error;
    \item the service-specific threshold tuner decides whether to emit an
    anomaly event;
    \item the coordinator batches related anomaly events and queries the RCA
    engine;
    \item the alert is enriched with likely root causes and propagation paths;
    \item in the background, edge nodes train locally and submit compressed,
    optionally DP-protected updates for FedAvg aggregation.
\end{enumerate}

\section{Data Flow}
\label{sec:design_dataflow}

\textbf{Metric ingestion.} Each agent reads from Prometheus or StatsD-compatible
receivers at 10\,s intervals into a fixed-capacity ring buffer. Missing values
are forward-filled for short gaps and masked for longer gaps.

\textbf{Inference.} Once the buffer contains $W=50$ observations, the agent
constructs a $(1,W,F)$ tensor and evaluates the LSTM-AE (Config~C: $h{=}128$,
$L{=}2$).

\textbf{Alert creation.} Alert payloads include the service, timestamp,
reconstruction error, threshold, model version, and trace correlation ID.
Events within the coordinator's 60\,s window are batched for RCA.

\textbf{Federated learning.} Rounds occur every 300\,s. Agents train for five
local epochs, apply DP and compression, and submit updates via gRPC.

\section{Component Design}
\label{sec:design_components}

\subsection{Edge Agent}
\label{subsec:design_edge_agent}

The edge agent is a Python process with three cooperative loops:

\begin{itemize}
    \item \textbf{Inference loop}: reads metric windows, computes anomaly
    scores, applies the current threshold, and emits anomaly events;
    \item \textbf{Training loop}: downloads the latest global model, performs
    local training, applies DP and compression, and uploads the update;
    \item \textbf{Threshold loop}: every 10 observations, updates the
    \texttt{ThresholdTuner} with the latest detection feedback.
\end{itemize}

Two node profiles are supported. The \emph{lightweight} profile is stateless and
is intended for minimal containers or devices with limited local storage. It
joins a round, trains, submits an update, and discards transient state. The
\emph{standard} profile persists the last global model, pending updates, Q-table
state, and local confusion counters in SQLite, allowing recovery after a crash
without restarting threshold learning from zero.

\subsection{Central Coordinator}
\label{subsec:design_coordinator}

The coordinator exposes the following gRPC endpoints:

\begin{itemize}
    \item \texttt{RegisterNode}: registers a node profile, software version,
    hardware capacity, and sample count;
    \item \texttt{SubmitUpdate}: receives compressed model updates and stores
    them in the current round buffer;
    \item \texttt{GetGlobalModel}: returns the latest model checkpoint and
    compression metadata;
    \item \texttt{ReportAnomaly}: accepts anomaly events for RCA batching.
\end{itemize}

Coordinator state is kept in SQLite with Write-Ahead Logging (WAL). The database
contains model checkpoints, round metadata, node registry entries, causal graph
snapshots, and alert records. SQLite is sufficient for the evaluated
single-coordinator deployment and avoids introducing an additional database
service. The design still permits migration to PostgreSQL or a replicated store
if multi-coordinator deployments are required.

Table~\ref{tab:persistence_schema} documents the minimum state that must survive
a restart for the framework to remain reproducible and auditable. It is not
intended as a complete physical database schema; instead, it identifies the
logical records required to reconnect experimental results to the model,
participants, graph snapshot, and alert context that produced them.

\begin{table}[H]
    \centering
    \small
    \begin{tabular}{|>{\raggedright\arraybackslash}p{0.24\textwidth}|
                    >{\raggedright\arraybackslash}p{0.29\textwidth}|
                    >{\raggedright\arraybackslash}p{0.35\textwidth}|}
        \hline
        \textbf{Store} & \textbf{Main fields} & \textbf{Purpose} \\
        \hline
        \texttt{model\_registry} & round, checksum, path, created\_at &
        versioned global checkpoints and rollback \\
        \hline
        \texttt{round\_log} & round, participants, status, duration &
        crash recovery and reproducible FL auditing \\
        \hline
        \texttt{node\_registry} & node id, profile, last seen, sample count &
        participation tracking and dropout handling \\
        \hline
        \texttt{causal\_graph} & snapshot id, nodes, edges, weights &
        RCA continuity across restarts \\
        \hline
        \texttt{alerts} & service, score, threshold, RCA candidates &
        incident history and threshold feedback \\
        \hline
    \end{tabular}
    \caption{Coordinator persistence records required for recovery and auditing.}
    \label{tab:persistence_schema}
\end{table}

These records define failure boundaries: the \texttt{round\_log} explains which
updates entered a global model, the \texttt{model\_registry} allows rollback,
and the \texttt{alerts} table preserves threshold and RCA candidates at
incident time.

\subsection{Trace Collector and RCA Engine}
\label{subsec:design_trace_collector}

The trace collector ingests spans from OpenTelemetry or Jaeger exporters,
extracts parent--child service relationships, and updates the
\texttt{CausalGraph} stored in NetworkX. Weak edges are pruned periodically, and
snapshots are persisted every 300\,s. When the RCA engine receives a batched
anomalous set, it computes candidate scores and returns a ranked root-cause
list. The RCA result is stored with the alert to preserve auditability, since
the causal graph evolves over time and recomputation could produce a different
explanation.

\section{Configuration and Deployment}
\label{sec:design_config}

Runtime configuration is supplied through environment variables and a
\texttt{config.yaml} file specifying model architecture ($h{=}128$, $L{=}2$,
dropout~0.2), federated settings (5~local epochs, minimum 3~clients,
120\,s timeout), privacy parameters ($\varepsilon{=}10$, $\delta{=}10^{-5}$,
$C{=}1.0$), compression algorithm selection (Zstd default, LZ4 above 80\% CPU),
and Q-learning hyperparameters ($\eta{=}0.1$, $\gamma{=}0.95$,
$\varepsilon_{\text{explore}}{=}0.1$). Non-destructive parameters can be
reloaded without restarting the coordinator.

For Kubernetes deployment, edge agents may be installed as a DaemonSet or as
sidecars. The coordinator runs as a Deployment backed by a persistent volume.
Configuration is provided through ConfigMaps, while TLS certificates and
webhook tokens are stored as Secrets.

\textbf{Security boundaries.} The implementation assumes an honest-but-curious
coordinator for the DP experiments. gRPC endpoints run over mutual TLS, node
registration tokens are scoped per deployment, and the coordinator validates
model update metadata (version, tensor shape, checksum) before aggregation.
These checks prevent accidental corruption but do not replace Byzantine robust
aggregation.

\textbf{Operational dashboards.} Grafana panels track anomaly scores per
service, FL round status, causal graph health, and communication overhead.

\section{Reliability and Fault Tolerance}
\label{sec:reliability_overview}

The framework prioritizes availability and partition tolerance. Edge agents keep
the last valid global model and continue local inference during coordinator
outages. FL rounds are allowed to proceed with partial participation when the
minimum client threshold is met. RCA may operate on a temporarily stale causal
graph during trace interruptions, but graph state converges as buffered spans are
replayed.

The main single points of failure and mitigations are:

\begin{itemize}
    \item \textbf{Coordinator failure}: model checkpoints and round state are
    written before commit, enabling restart from the last valid checkpoint;
    \item \textbf{Trace collector failure}: spans are buffered before graph
    ingestion, and periodic graph snapshots prevent full topology loss;
    \item \textbf{Edge agent failure}: lightweight agents re-register, while
    standard agents replay pending updates and reload their cached Q-table;
    \item \textbf{Network partition}: agents continue inference locally and
    rejoin federated training when connectivity returns.
\end{itemize}

Heartbeat monitoring marks nodes offline after a configurable timeout; once a
heartbeat resumes, they rejoin the next round.

Table~\ref{tab:recovery_behavior} shows recovery behavior for each state
category.

\begin{table}[H]
    \centering
    \small
    \begin{tabular}{|p{0.22\textwidth}|p{0.30\textwidth}|p{0.36\textwidth}|}
        \hline
        \textbf{State} & \textbf{Failure symptom} & \textbf{Recovery action} \\
        \hline
        Inference window & scrape gap or local restart & reuse cached model;
        mask long metric gaps \\
        \hline
        Pending update & upload interrupted & replay from standard profile
        cache or discard in lightweight mode \\
        \hline
        Round aggregate & coordinator restart & reload last committed round;
        abort incomplete aggregate \\
        \hline
        Causal graph & trace collector outage & replay buffered spans and
        resume from latest graph snapshot \\
        \hline
        Threshold policy & agent crash & reload Q-table in standard profile;
        otherwise restart from safe threshold \\
        \hline
    \end{tabular}
    \caption{Recovery behavior for critical framework state}
    \label{tab:recovery_behavior}
\end{table}

This model favors graceful degradation: during partitions, local detection may
use a stale model and RCA a stale graph, which are acceptable for monitoring
since stopping detection entirely would be worse.

\section{Implementation Stack}
\label{subsec:design_stack}

The implementation uses Python~3.12, PyTorch~2.x, gRPC with Protocol Buffers,
Zstandard, LZ4, NetworkX, SQLite in WAL mode, Prometheus-compatible metric
clients, the OpenTelemetry SDK, Docker, and Kubernetes. The stack was selected
for reproducibility and low operational overhead. PyTorch keeps the LSTM-AE and
DP gradient manipulation transparent. gRPC gives typed service contracts and
efficient binary transport. SQLite with WAL offers durable state without
requiring a separate database cluster, which is appropriate for the
single-coordinator system evaluated in this thesis.

The system design described in this chapter translates the algorithmic
methodology of Chapter~\ref{ch:framework} into deployable components with
explicit failure boundaries, persistence contracts, and configuration surfaces.
Chapter~\ref{ch:evaluation} now evaluates this implementation empirically:
detection accuracy, privacy--utility trade-offs, communication overhead, RCA
quality, adaptive thresholding behavior, and coordinator scalability.
